\scanpage{33}
\mspage{15}

\textcircled{2} \uline{Le groupe} $S\widehat{\widehat{\unsure{\Sigma|E}}}$.

Par le calcul, il s'explicite comme l'ens des systèmes

\[
(\underset{\substack{\rotatebox{-90}{$\in$}\\\widehat{\mathbb Z}^{*}}}{p},
 \underset{\substack{\rotatebox{-90}{$\in$}\\\unsure{\mathfrak S_3|\mathcal S_3}}}{\sigma},
 \underset{\substack{\rotatebox{-90}{$\in$}\\\widehat\pi}}{g_0},
 \underset{\substack{\rotatebox{-90}{$\in$}\\\widehat\pi}}{g_1},
 \underset{\substack{\rotatebox{-90}{$\in$}\\\widehat\pi}}{g_\infty})=\widetilde u \tag*{(1)}
\]

satisfaisant la condition

\[
\mathrm{int}\unsure{(g)_{\infty}|(g_{\infty})}(l_{\sigma(\infty)}^p)\cdot
\mathrm{int}(g_1)(l_{\sigma(1)}^p)\cdot
\mathrm{int}(g_\infty)(l_{\sigma(\unsure{\infty|0})}^p)=1 \tag*{(2)}
\]

et tels que les trois facteurs de ce produit soient un syst. de générateurs de $\widehat\pi$ [avec la relation génératrice (2)]. Le morphisme $u$ défini par ce \deleted{\unsure{morphi}} système est

\[
\left\{\begin{aligned}
u(l_0)&\unsure{\mapsto|=}\mathrm{int}(g_0)\,l_{\sigma(0)}^p\\
u(l_1)&=\mathrm{int}(g_1)\,l_{\sigma(1)}^p\\
u(l_\infty)&=\mathrm{int}(g_\infty)\,l_{\sigma(\infty)}^p
\end{aligned}\right. \tag*{(3)}
\]

ce qui explicite l'homom

\[
S\widehat{\widehat\Sigma}\xrightarrow{p}\widehat{\widehat\Sigma}=\mathrm{Aut}_{\mathrm{lac}}(\widehat\pi) \tag*{(4)}
\]

L'homomorphisme canonique

\[
i:\widehat\pi\times\widehat{\mathbb Z}^{\{0,1,\infty\}}\longrightarrow S\widehat{\widehat\Sigma} \tag*{(\unsure{5})}
\]

s'explicite par

% slanted
\margin{On \unsure{retrouve|retiendra} que (\unsure{4|p}) est surjectif, et a comme noyau $i(\widehat\pi)$}

\[
i(g;\alpha_0,\alpha_1,\alpha_\infty)=(1,1,g\,l_0^{\alpha_0},g\,l_{\unsure{0|1}}^{\alpha_1},g\,l_\infty^{\alpha_\infty})\quad\text{donc} \tag*{(6)}
\]

\[
\deleted{\text{donc}}\quad pi(g;\alpha_0,\alpha_1,\alpha_\infty)=\mathrm{int}\,g. \tag*{(7)}
\]

La multiplication dans $\widehat{\widehat\Sigma}$ est \unsure{donné|donnée} par

\[
\underbrace{(p,\sigma,g_0,g_1,g_\infty)}_{\widetilde u}
\underbrace{(p',\sigma',g'_0,g'_1,g'_\infty)}_{\widetilde u'}
=(pp',\sigma\sigma',u(g'_0)g_{\sigma'(0)},u(g'_1)g_{\sigma'(1)},u(g'_\infty)g_{\sigma'(\infty)}\unsure{)} \tag*{(8)}
\]

Cette formule n'est pas, du point de vue du calcul,
% archive sheet number: 32
